% SOURCE_STATEMENT: "For the Pauli–Fierz Hamiltonian $H(p)$ with photon dispersion $\\omega(k)=\\sqrt{|k|^2+m_{\\mathrm{ph}}^2}$, does its $g$-factor have a finite limit as $m_{\\mathrm{ph}}\\to0$ and $\\widehat\\varphi(k)\\to(2\\pi)^{-3/2}$? If the point-charge limit exists, determine its dependence on the fine-structure constant $\\alpha$, including the coefficient of $O(\\alpha^2)$ in $$g_\\infty=2\\left(1+\\frac{4\\alpha}{3\\pi}+O(\\alpha^2)\\right)$$."
% Preparation and final self-audit: CONTRIBUTING.md prompts C and D; source/audit artifacts retained in the local physics_statements_2026_10_11 batch.
% !TeX program = lualatex
\documentclass[11pt]{article}
\usepackage[margin=25mm]{geometry}
\usepackage{fontspec}
\setmainfont{Latin Modern Roman}
\usepackage{amsmath,amssymb,mathtools}
\usepackage{unicode-math}
\setmathfont{Latin Modern Math}
\usepackage{xurl,hyperref}
\hypersetup{colorlinks=true,urlcolor=blue}
\setcounter{secnumdepth}{0}
\setlength{\parindent}{0pt}
\setlength{\parskip}{5pt}
\setlength{\emergencystretch}{3em}
\newcommand{\sref}[2]{\hyperlink{source:#1:#2}{[S#2]}}
\newcommand{\cataloguescope}[1]{\paragraph{Recorded target.}#1}
\newcommand{\R}{\mathbb R}
\newcommand{\C}{\mathbb C}
\newcommand{\N}{\mathbb N}
\newcommand{\Z}{\mathbb Z}
\newcommand{\Q}{\mathbb Q}
\newcommand{\F}{\mathbb F}
\newcommand{\eps}{\varepsilon}
\begin{document}
% UnsolvedMath ID 30000925; source catalogue code OWR-1789-002
\setcounter{section}{30000924}
\section{Massless and Point-Charge Limits of the Pauli–Fierz $g$-Factor}\label{30000925}
{\small UnsolvedMath ID 30000925 · Mathematical Physics}
\subsection{Definitions and mathematical statement}

\paragraph{Shared notation.}$\mathbb N=\{1,2,\ldots\}$; the standard number fields have their usual meanings. All additional conventions are stated locally.


\paragraph{Definitions and conventions.} The translation-invariant Pauli--Fierz Hamiltonian decomposes into total-momentum fibers $H(p)$ acting on electron-spin and photon Fock space. $m_{\mathrm{ph}}$ regularizes the photon infrared dispersion and $\widehat\varphi(k)$ is the charge form factor providing the ultraviolet cutoff. The $g$-factor measures the spin response relative to the effective orbital magnetic response and is defined by the report's equation (14). The point-charge limit sends $\varphi$ to a delta distribution, equivalently $\widehat\varphi(k)\to(2\pi)^{-3/2}$ in its Fourier convention. The fine-structure constant is $\alpha=e^2/(4\pi)$ in the report's units.

\paragraph{Source scope and qualifications.} The perturbative first correction is not a proof that the full cutoff limits exist. The two cutoff removals, possible order of limits, all-alpha dependence, and the second-order coefficient are separate questions. The g-factor definition is the source's equation (14), including effective-mass normalization.

\paragraph{Catalogue mathematical target.}
\paragraph{Catalogue wording retained for traceability.}
For the Pauli–Fierz Hamiltonian $H(p)$ with photon dispersion $\omega(k)=\sqrt{|k|^2+m_{\mathrm{ph}}^2}$, does its $g$-factor have a finite limit as $m_{\mathrm{ph}}\to0$ and $\widehat\varphi(k)\to(2\pi)^{-3/2}$? If the point-charge limit exists, determine its dependence on the fine-structure constant $\alpha$, including the coefficient of $O(\alpha^2)$ in $$g_\infty=2\left(1+\frac{4\alpha}{3\pi}+O(\alpha^2)\right)$$.

\paragraph{Saved literature context.} The supplied catalogue's assessment dated 2026-08-21 records: Perturbative coefficients and several ultraviolet or point-limit results are known, but the report's full infrared limit, alpha-dependence, and O(alpha squared) questions for the interacting Hamiltonian remain unresolved. The assessment is retained as dated provenance; it is not a new comprehensive literature or open-status certification.
\subsection{Short English statement}
Determine whether the Pauli--Fierz g-factor survives removal of both photon-mass and charge-size cutoffs, and compute its dependence on coupling.
\subsection{Sources}
\begin{itemize}
\item[\hypertarget{source:30000925:1}{S1}] ULAM.AI and the UnsolvedMath Contributors, UnsolvedMath, record 30000925 (OWR-1789-002). Catalogue curation under CC BY 4.0; definitions and English exposition added locally.
\url{https://huggingface.co/datasets/ulamai/UnsolvedMath}.
\item[\hypertarget{source:30000925:2}{S2}] Mini-Workshop: Attraction to Solitary Waves and Related Aspects of Physics (2008), Oberwolfach Reports, DOI 10.4171/owr/2008/08, pp. 367--400; official report source and mathematical context.
\url{https://doi.org/10.4171/owr/2008/08}.
\item[\hypertarget{source:30000925:3}{S3}] Official EMS Press PDF of the same Oberwolfach report; PDF page 29, printed page 395, equations (14)--(16).
\url{https://ems.press/content/serial-article-files/46151}.
\item[\hypertarget{source:30000925:4}{S4}] Pauli--Fierz point-limit dynamics; reference in the saved dated catalogue literature review.
\url{https://arxiv.org/abs/math-ph/0407002}.
\item[\hypertarget{source:30000925:5}{S5}] ultraviolet-limit analysis; reference in the saved dated catalogue literature review.
\url{https://arxiv.org/abs/2004.06494}.
\end{itemize}
\label{end:30000925}
\end{document}
